% 06_discussion -- written by Claude Code (D-045); revised after the peer review (Prompt 16: M1, M3, M4, M5, m20).
% Mechanisms are conjectures and are worded as such.
\section{Discussion}\label{sec:discussion}

\paragraph{A modest, uncertain gain, mostly from voting.} The majority of 10 sampled answers was more accurate than a
single greedy answer; the gain is statistically supported and its point estimate passes the threshold set in advance,
but the data are compatible with any gain from \diffMainLo{} to \diffMainHi{} points, it costs \tokRatio{} times the
output tokens, and most of it comes from aggregating runs rather than from sampling itself. Self-consistency was
proposed on the intuition that several reasoning paths converge on the correct answer more often than a single
one~\cite{wang2022selfconsistency}; our data fit this view, since greedy's wrong answers were rare among the sampled
runs, but do not prove a mechanism. A second greedy run, parse failures, ties and the MedQA split of origin do not
change the direction. Strong closed models gained little from voting in
oncology~\cite{holmes2023,stifini2026,windisch2026}; a model that is already consistent has less to gain, which may
explain the difference, but with one model tested this remains a hypothesis.

\paragraph{Consistency depends on the model and its decoding.} Greedy decoding repeated its answer for
\rerunSamePct\% of items, two sampled runs agreed for \pairAgree\%, and unanimity fell between $T$ = 0.3 and 1.0. That
consistency varies with temperature was shown before for small open models~\cite{pinhanez2025} and clinical text
extraction~\cite{liu2026repro}, and that it varies between models in a tumor-board study~\cite{li2025imitation}; we add
the same picture for reasoning outputs on oncology questions. Our figures cannot be ranked against the ChatGPT
tumor-board studies~\cite{liao2025,asoglu2026} (real cases, free text, three runs, unreported settings), and our 9-of-10
stability (\stableTHigh--\stableTLow\%) is only loosely comparable with letter-only answers~\cite{pinhanez2025}. The
number of runs and their consistency, which our review recommended reporting~\cite{sharma2026review}, can only be
interpreted together with the decoding settings.

\paragraph{Agreement is a weak signal.} Agreement tracked correctness (AUROC \auroc), but it needs many runs, partly
reflects item difficulty, and \unanimousWrongPct\% of unanimous answers were wrong, all of them also wrong under greedy
decoding. Such systematic errors, presumably gaps in knowledge or reasoning, and the runs' shared preference against
option A are invisible to both voting and agreement.

\paragraph{Implications.} These results concern exam questions with a single keyed answer, not board decisions, and say
nothing about patient care or clinical use of model output. For studies that compare LLMs with tumor boards they
support one reporting practice: state how many times each case was run, with which decoding settings, and how
consistent the runs were.
